% Sommaire et analyse de la réponse de l'arrondissement à la demande d'accès
% 2026-581 (étude de la rue de Terrebonne). Version française.
% Compiler avec pdflatex (deux fois).
\documentclass[11pt,letterpaper]{article}
\usepackage[T1]{fontenc}
\usepackage[utf8]{inputenc}
% babel-french n'est pas disponible ici : typographie française appliquée à la main
% (espaces insécables avant : ; et à l'intérieur des guillemets).
\usepackage[scaled=0.95]{helvet}
\renewcommand{\familydefault}{\sfdefault}
\usepackage[letterpaper,top=0.9in,bottom=0.95in,left=1in,right=1in]{geometry}
\usepackage{microtype}
\usepackage{enumitem}
\usepackage{titlesec}
\usepackage{fancyhdr}
\usepackage{xcolor}
\usepackage{booktabs}
\usepackage{longtable}
\usepackage{array}
\usepackage[hidelinks]{hyperref}
\urlstyle{same}
\newcommand{\docversion}{2026-10-13-15-12-54} % set by build.py (America/Toronto)
\hypersetup{pdfsubject={version \docversion}}

\setlength{\parindent}{0pt}
\setlength{\parskip}{0.5em}
\linespread{1.03}

\titleformat{\section}{\normalfont\large\bfseries}{\thesection}{0.6em}{}[\vspace{1pt}\titlerule]
\titlespacing*{\section}{0pt}{1.1em}{0.4em}
\titleformat{\subsection}{\normalfont\bfseries}{\thesubsection}{0.6em}{}
\titlespacing*{\subsection}{0pt}{0.8em}{0.25em}

\pagestyle{fancy}
\fancyhf{}
\renewcommand{\headrulewidth}{0pt}
\fancyfoot[C]{\footnotesize\color{gray}Demande d'accès 2026-581 -- sommaire et analyse -- page \thepage}
\fancyfoot[R]{\tiny\color{gray}v.\,\docversion}

% ---- Liens vers le dépôt ----
\newcommand{\repourl}{https://github.com/OpenMobilityData/Terrebonne-access-requests}
\newcommand{\repo}[1]{\href{\repourl}{#1}}
\newcommand{\doc}[2]{\href{\repourl/blob/main/#1}{#2}}
% GitHub Pages : les ancres #page=N fonctionnent ici, pas sur la visionneuse de GitHub.
\newcommand{\pagesurl}{https://openmobilitydata.github.io/Terrebonne-access-requests}
\expandafter\def\csname bundle01\endcsname{2026-581/raw/01_emails_mandate_counts.pdf}
\expandafter\def\csname bundle02\endcsname{2026-581/raw/02_framework_agreement_tender.pdf}
\newcommand{\pgl}[3]{\href{\pagesurl/\csname bundle#1\endcsname\#page=#2}{\textsf{\small[#1, p.~#3]}}}
\newcommand{\pg}[2]{\pgl{#1}{#2}{#2}}
\newcommand{\quot}[1]{\begin{quote}\small #1\end{quote}}

\begin{document}

\begin{center}
{\Large\bfseries Demande d'accès 2026-581~: ce que révèle le dossier de l'arrondissement}\\[0.4em]
{\large Sommaire et analyse des documents communiqués le 9~octobre 2026}\\[0.6em]
Richard Hoge, Côte-des-Neiges--Notre-Dame-de-Grâce -- 9~octobre 2026 (révisé le 11~octobre)\\
Publié avec les documents sources à \repo{\repourl}\\
{\small Traduction de la version anglaise\,; en cas de divergence, l'anglais fait foi.}
\end{center}

\section*{Sommaire}

Le 11~septembre 2026, j'ai demandé à l'arrondissement de Côte-des-Neiges--Notre-Dame-de-Grâce
les documents relatifs à l'étude d'ingénierie qu'il avait commandée sur la
rue de Terrebonne~: qui l'a demandée, pourquoi, ce qu'elle devait examiner,
ce qu'elle a coûté et ce qu'elle a livré. L'arrondissement a répondu le
9~octobre par 508~pages (dossier 2026-581). L'étude elle-même n'a pas été
communiquée. Ce qui l'a été suffit à établir ce qui suit.

\begin{enumerate}[nosep]
  \item \begin{minipage}[t]{\linewidth}\textbf{La décision a précédé l'étude.} Le 10~juin 2026, la
    division de l'ingénierie de l'arrondissement écrivait à la division de
    l'aménagement des rues de la Ville que les élus souhaitaient revoir
    l'aménagement de 2024 «\,avec l'objectif de redonner la bande de
    stationnement qui avait été retranchée\,», que les ingénieurs de
    l'arrondissement avaient \emph{déjà} émis un avis technique sur les deux
    concepts esquissés par les élus, et que ceux-ci voulaient «\,un avis
    officiel, neutre et externe\,». Le mandat finalement confié à EXP le
    20~août a pour titre, en partie, «\,évaluer la place du stationnement sur
    rue, dans le but d'en augmenter le nombre\,».\end{minipage}\vspace{0.35em}
\item \begin{minipage}[t]{\linewidth}\textbf{L'équipe vélo de la Ville elle-même a soulevé des enjeux
    de sécurité.} Le Service de l'urbanisme et de la mobilité a remis une
    analyse en mode express le 17~juin\,; le résumé qu'en fait
    l'arrondissement, rédigé cinq jours plus tard, indique que les options
    «\,présentent plusieurs compromis importants sur le plan de la sécurité,
    du confort des usagers vulnérables et de la cohérence du réseau
    cyclable\,». Cette analyse n'a pas été communiquée. L'avis du Service de
    sécurité incendie du 13~juillet est caviardé en entier.\end{minipage}\vspace{0.35em}
\item \begin{minipage}[t]{\linewidth}\textbf{La mesure hivernale est qualifiée de transitoire\,; le
    mandat du cabinet de la mairesse était permanent.} Le 13~août, la cheffe
    de cabinet de la mairesse d'arrondissement a fixé le mandat par écrit~:
    une piste cyclable bidirectionnelle d'un côté de Terrebonne, de Loyola à
    Côte-Saint-Luc, «\,qui respecte les normes minimales prévues par la loi\,»,
    et «\,le retour du stationnement d'un côté de la rue\,». Aucune mesure
    hivernale n'y figure. La demande de services, une semaine plus tard, a
    ajouté un Volet~1, le retrait des aménagements cyclables en hiver, à
    titre de «\,mesure transitoire pour l'hiver 2026--2027\,», avant le
    Volet~2, la piste bidirectionnelle permanente. La conversion hivernale
    annoncée le 6~octobre est la première étape d'une réduction à l'année de
    l'espace cyclable qui n'a jamais été divulguée au public.\end{minipage}\vspace{0.35em}
\item \begin{minipage}[t]{\linewidth}\textbf{L'«\,étude\,» est une note technique de six semaines et de
    14\,755~\$, avec une date de conseil fixée d'avance.} Demande de
    services le 20~août\,; offre le 25~août\,; bon de commande le 26~août\,;
    réunion de démarrage le 3~septembre\,; version préliminaire le
    17~septembre\,; version finale due le 2~octobre\,; annonce le 6~octobre. Le
    personnel de l'arrondissement a écrit au consultant le 18~août que
    «\,l'objectif de nos élus est de pouvoir agir sur vos conclusions pour un
    horizon début octobre\,» et, le 10~septembre, que la version préliminaire
    devait être livrée au plus tard le 18 «\,pour respecter notre calendrier
    du Conseil d'arrondissement du 5~octobre\,».\end{minipage}\vspace{0.35em}
\item \begin{minipage}[t]{\linewidth}\textbf{Les comptages de vélos cités au conseil reposent sur les
    données les plus faibles du dossier et sous-estiment l'usage
    hivernal\,; le dossier contenait les données du compteur permanent pour
    l'hiver dernier, qui attestent un usage hivernal moyen au moins cinq fois
    plus élevé que le chiffre cité. L'annonce du conseiller n'a pas tenu compte
    de ces données plus robustes.} Le conseiller a dit à la séance (en
    anglais) que «\,we have counters everywhere, we have more counters than
    before. So we know exactly how many people use it, when and how\,» et
    que «\,we go from 700 during the good months to 8 \ldots\ So that's 4
    cyclists, right? Including me\,». Le «\,8\,», et le «\,12\,» qu'il a cité
    en réponse à une question du public, proviennent de la
    présentation des comptages de l'arrondissement de février 2025~: il ne
    s'agit ni de vélos par jour ni de vélos par heure, mais de vélos comptés
    au moyen d'une caméra temporaire (Miovision) installée aux intersections
    Cavendish et Grand pendant six heures de pointe d'une seule journée, le
    29~janvier 2025, lors du premier hiver de la piste, alors que le
    déneigement était souvent en retard, et \emph{une journée où Environnement
    Canada avait émis un avis de temps hivernal et où la population était
    invitée à reporter les déplacements non essentiels}\,; la même caméra a
    enregistré une baisse d'environ le quart des automobiles sur les approches
    Terrebonne et des piétons par rapport au comptage de novembre, et de 95~\%
    des vélos. La diapositive ne précise rien de tout cela. Avant le premier courriel du dossier (10~juin) et le
    mandat du cabinet du maire (13~août), un compteur permanent avait déjà
    enregistré chaque vélo passé sur la rue pendant tout l'hiver suivant,
    soit une année entière plus récente\,; les ingénieurs de l'arrondissement
    ont offert ses données à la Ville le 10~juin, et son tableau était entre
    les mains du consultant en août~: 4~891 passages de décembre à mars (41
    par jour\,; 46 les jours de semaine), et un deuxième été à environ 700
    par jour. Pour la question qui touche à la sécurité, l'exposition, le
    chiffre pertinent est le nombre total de déplacements hivernaux, et non
    le comptage d'une journée choisie\,; ce total était en la possession de
    l'arrondissement et n'a pas été cité au conseil. La baisse
    de «\,99~\%\,» citée à la séance, 700 contre 8, oppose le plus grand nombre
    d'une diapositive, Grand en septembre 2024 compté sur les quatre
    approches, au plus petit nombre de la suivante, Grand en janvier 2025 sur
    l'axe Terrebonne seulement\,; comparées de façon équivalente, les
    diapositives elles-mêmes donnent 97 à 98~\%, étiquettent janvier à
    $-95$~\%, et le compteur donne environ 93~\% d'une saison entière à
    l'autre. Les deux sources s'entendent sur le fait que la baisse est
    forte\,; elles diffèrent sur le nombre absolu, et une seule d'entre elles
    donne «\,quatre personnes\,». Quelle qu'en soit l'explication, une mesure
    présentée comme fondée sur les données reposait sur six heures
    de comptage en une seule journée et a ignoré des milliers d'heures de
    comptage plus récent sur tout un hiver.\end{minipage}\vspace{0.35em}
\item \begin{minipage}[t]{\linewidth}\textbf{Le besoin de stationnement hivernal supplémentaire n'a jamais été mesuré.}
    Dans les 508~pages, il n'y a aucun inventaire des cases sur rue, aucun
    relevé d'occupation, aucun sondage auprès des résidents et aucune
    estimation du nombre de cases que la conversion créerait\,; les «\,50 à
    60 cases\,» que le conseiller a données à l'APC-NDG lors de la rencontre
    de consultation ne figurent nulle part dans les documents
    communiqués. Le mandat demande au consultant de compter les cases qu'un
    aménagement «\,pourrait permettre de recréer\,» et de mesurer
    «\,l'amélioration de la disponibilité du stationnement\,»\,; l'offre
    promet une analyse de «\,l'impact sur le nombre de places\,» et ne
    prévoit aucun travail de terrain\,; le consultant a travaillé à partir
    d'imagerie aérienne. Ce que la note dit du stationnement est donc un
    décompte de cases qu'on pourrait tracer, non une mesure du besoin. Rien
    dans le dossier n'indique combien de nuits d'hiver le stationnement
    près de Terrebonne est réellement complet.\end{minipage}\vspace{0.35em}
\item \begin{minipage}[t]{\linewidth}\textbf{L'argument en faveur du retrait d'une piste repose sur l'idée
    que les pistes sont plus larges que leur circulation cycliste ne le
    justifie\,; cette idée est l'opinion des élus, non l'avis des
    ingénieurs.} Le courriel du 10~juin rapporte la préoccupation des élus
    selon laquelle la largeur des pistes était «\,jugée disproportionnée par
    rapport au volume de cyclistes observé \ldots\ selon les élus\,», et le
    conseiller Teodoresco a repris le point devant l'APC-NDG lors de la
    rencontre de consultation, mais non à la séance du conseil du
    6~octobre. Aucun ingénieur au dossier ne l'endosse. Les ingénieurs de
    l'arrondissement décrivent les solutions de rechange comme «\,une voie de
    circulation très étroite (moins de 4~mètres)\,» et «\,un aménagement
    bidirectionnel très étroit\,»\,; l'équipe vélo de la Ville y a vu
    «\,plusieurs compromis importants\,»\,; et le mandat du 13~août demande
    un aménagement «\,qui respecte les normes minimales prévues par la
    loi\,», ce qui revient à admettre que les concepts de juin ne les
    respectaient pas. La largeur d'une piste est fixée par le gabarit du
    cycliste, l'espace de dépassement et le matériel de déneigement, dont
    aucun ne dépend du nombre de personnes qui roulent\,; rien au dossier ne
    décrit les pistes actuelles comme plus larges que la norme ne
    l'exige.\end{minipage}\vspace{0.35em}
\item \begin{minipage}[t]{\linewidth}\textbf{Une option limitée aux opérations de déneigement
    faisait partie du mandat.} L'offre d'EXP énumère, parmi les scénarios à
    comparer, «\,la fermeture en période de déneigement seulement\,». Le
    conseiller a dit à l'APC-NDG, lors de la rencontre qu'il a ensuite citée
    au conseil comme sa consultation de l'Association, que le stationnement
    limité aux tempêtes avait été écarté parce qu'on ne pouvait en informer
    les résidents\,; l'option n'a pas été mentionnée dans son allocution du
    6~octobre, et l'analyse sur laquelle repose cette mise à l'écart se
    trouve dans la note qui n'a pas été communiquée.\end{minipage}\vspace{0.35em}
\item \begin{minipage}[t]{\linewidth}\textbf{L'exercice ne suit pas la méthode que la Ville a elle-même
    publiée pour réviser une piste cyclable.} Le 10~février 2026,
    l'administration a lancé son audit du réseau cyclable sur le principe
    «\,recueillir les données, analyser et corriger au besoin\,»~: analyse
    technique, collecte de données et observation terrain pendant la haute
    saison, sous la direction du Service de l'urbanisme et de la mobilité,
    avec des résultats attendus à la fin de 2026\,; le même communiqué nomme
    Terrebonne parmi les rues où la Ville travaille avec l'arrondissement sur
    la sécurité. Sur Terrebonne, l'objectif a été fixé avant toute analyse,
    l'analyse du SUM a été défavorable et refusée, le mandat visait à
    augmenter le stationnement, aucune donnée n'a été recueillie, et le
    changement a été annoncé alors que l'audit de la Ville était toujours en
    cours.\end{minipage}\vspace{0.35em}
\item \begin{minipage}[t]{\linewidth}\textbf{Le refus de communiquer le rapport technique repose sur un
    motif qui s'éteint avec la décision.} La note et l'analyse de la Ville sont
    refusées en vertu de l'article~39, qui s'applique «\,jusqu'à ce que la
    recommandation ait fait l'objet d'une décision\,». Le 6~octobre, le
    conseiller a dit que les élus étaient «\,now ready to give our
    orientation to the services\,» et que le retrait saisonnier «\,is what
    we're going to propose\,»\,; sa déclaration signée ajoute que «\,one
    thing is certain: residents will see changes before Christmas 2026\,».
    Ou bien la recommandation a fait l'objet d'une décision et le motif
    n'existe plus, ou bien elle ne l'a pas été et le résultat a été annoncé
    avant la décision. Une demande de révision peut être
    adressée à la Commission d'accès à l'information dans les 30~jours, soit
    d'ici le \textbf{8~novembre 2026}.\end{minipage}\vspace{0.35em}
\end{enumerate}

\section{Ce qui a été demandé et ce qui a été reçu}

La demande (11~septembre 2026\,; accusé de réception le 14~septembre sous le
numéro 2026-581) comportait cinq volets~: A, origine et objet\,; B,
approvisionnement et contrat\,; C, méthodologie et exécution\,; D, livrables
et révision\,; E, traitement interne et communications. La lettre de réponse
traite chaque point l'un après l'autre et joint huit fichiers. Les deux
liasses substantielles sont~:

\begin{itemize}[nosep]
  \item \textbf{Liasse 01} (\doc{2026-581/raw/01_emails_mandate_counts.pdf}{01 Doc\_req
    2026-581.pdf}, 167~pages)~: les chaînes de courriels, la demande de
    services, l'offre de services d'EXP, le bon de commande, le formulaire de
    la réunion de démarrage, la présentation d'analyse des comptages de
    février 2025 et sa mise à jour de 2025, les totaux quotidiens de la
    boucle de comptage Kensington (du 26~juillet 2025 au 31~août 2026) et les
    tableaux de comptage Miovision aux intersections pour 2024--2025.
  \item \textbf{Liasse 02} (\doc{2026-581/raw/02_framework_agreement_tender.pdf}{02 DOC requis
    2026-581.pdf}, 341~pages)~: l'appel d'offres public pour les
    ententes-cadres 2025--2027 (59-25-ET-01), le règlement de gestion
    contractuelle, le dossier décisionnel et les résultats du comité de
    sélection, ainsi que la soumission gagnante de WSP pour le contrat~1.
    Rien n'y concerne Terrebonne en particulier\,; la liasse documente le
    véhicule contractuel par lequel EXP a été retenue.
\end{itemize}

Les fichiers communiqués sont publiés tels que reçus dans le dépôt
\repo{OpenMobilityData/Terrebonne-access-requests} sur GitHub, sous
\texttt{2026-581/raw/}\,; le matériel dérivé, dont le présent document, est
tenu à part sous \texttt{2026-581/derived/}. Les autres pièces jointes sont
la résolution du conseil octroyant les ententes-cadres (CA25~170126,
9~juin 2025), l'avis de recours et des copies des articles de la Loi
invoqués (9, 23, 24, 37, 39, 53 à 59). La lettre de réponse se trouve à
\doc{2026-581/raw/response_2026-10-09.pdf}{Request 2026-581}\,; la demande
originale à \doc{2026-581/raw/request_2026-09-11.pdf}{Terrebonne study
access request}.

\textbf{Points sans documents.} L'arrondissement n'a rien trouvé pour A2
(financement externe), C3 (réunions de suivi), E2 (plans de communication)
et E3 (documents de conception ou de mise en œuvre antérieurs à la
réception du rapport final). Le point C3 se concilie mal avec les
courriels communiqués, qui fixent les rencontres de travail n\textsuperscript{o}~1 et
n\textsuperscript{o}~2 sur Teams et renvoient aux notes qu'en a prises le consultant
\pgl{01}{23}{23--24, 167}\,; le formulaire de démarrage indique qu'«\,on
pourra produire des comptes rendus distincts pour les prochaines
rencontres\,». Ou bien ces comptes rendus n'ont pas été produits, ou bien on
ne les a pas cherchés.

\section{Chronologie établie par le dossier}

\begin{small}
\begin{longtable}{>{\raggedright\arraybackslash}p{2.6cm}>{\raggedright\arraybackslash}p{10.2cm}>{\raggedright\arraybackslash}p{2.0cm}}
\toprule
Date & Événement & Source \\
\midrule
\endhead
2022--2023 & Étude de circulation d'EXP pour la piste protégée (mandat précédent)\,; scénario retenu~: sens unique divergeant à partir de Cavendish, retrait d'une bande de stationnement. & \pg{01}{8} \\
2024 & Réaménagement réalisé\,; 2,5~km à sens unique, pistes unidirectionnelles des deux côtés. & \pg{01}{3} \\
Juil.\ 2024 -- juin 2025 & Comptages Miovision de l'arrondissement à Cavendish, Grand et Walkley\,; Telraam à King-Edward et Royal. & \pgl{01}{82}{82--158} \\
29 janv.\ 2025 & Comptages Miovision d'une seule journée à Cavendish et Grand~: 18 et 12 vélos sur l'axe Terrebonne en 12~heures\,; 10 et 8 pendant les deux périodes de pointe. & \pgl{01}{85}{85, 155} \\
Févr.\ 2025 & Présentation «\,Suivi de projet -- Analyse des comptages\,»\,; mise à jour plus tard en 2025 avec les comptages de janvier et de juin et les chiffres de $-95$/$-96$~\%. & \pgl{01}{27}{27--73} \\
26 juil.\ 2025 & Mise en service de la boucle de comptage Kensington. & \pg{01}{74} \\
Nov.\ 2025 & Élections municipales. & --- \\
10 juin 2026 & Le chef de division (Trottier) au SUM~: les élus veulent rétablir la bande de stationnement\,; deux coupes de rue dessinées (chaussée désignée\,; piste bidirectionnelle très étroite)\,; l'arrondissement a déjà émis un avis technique\,; les élus en veulent un externe\,; rencontre avec la mairesse d'arrondissement prévue le 18~juin. & \pgl{01}{7}{7--8} \\
17 juin 2026 & Le SUM remet son analyse technique express (non communiquée). & \pgl{01}{6}{6--7} \\
22 juin 2026 & L'arrondissement demande l'avis du SIM\,; résume la conclusion du SUM comme «\,plusieurs compromis importants\,» sur le plan de la sécurité. & \pgl{01}{5}{5--6} \\
13 juil.\ 2026 & Réponse du SIM (corps du texte caviardé). & \pg{01}{5} \\
13 août 2026 & La cheffe de cabinet (D'Amours) au directeur de l'aménagement urbain (Admo), c.c.\ le directeur d'arrondissement, la mairesse d'arrondissement et le conseiller Teodoresco~: la mairesse veut une étude externe formelle\,; mandat~: piste bidirectionnelle d'un côté, de Loyola à Côte-Saint-Luc, respectant «\,les normes minimales prévues par la loi\,», et retour du stationnement d'un côté. Transmis à la division de l'ingénierie le jour même. & \pgl{01}{1}{1--2} \\
18 août 2026 & Soulière à EXP~: la nouvelle mairesse d'arrondissement veut une note technique\,; la rapidité est ce qui compte pour les élus\,; l'objectif est d'agir à un conseil d'arrondissement au début d'octobre. & \pgl{01}{21}{21--22} \\
20 août 2026 & Demande de services 59-26-ET-01\_002 à EXP. & \pgl{01}{3}{3--4} \\
25 août 2026 & Offre d'EXP~: 14\,755~\$\,; note finale le 2~octobre. & \pgl{01}{9}{9--13} \\
26 août 2026 & Bon de commande 1796945, 17\,706~\$ avec contingence. & \pgl{01}{15}{15--17} \\
28 août 2026 & L'arrondissement envoie à EXP la présentation des comptages de 2025 et le tableau de la boucle de comptage\,; EXP envoie sa première analyse intermédiaire (côté sud privilégié pour une piste bidirectionnelle). & \pg{01}{19} \\
3 sept.\ 2026 & Réunion de démarrage (Teams). & \pg{01}{26} \\
10 sept.\ 2026 & Arrondissement~: version préliminaire requise pour le 18~septembre «\,pour respecter notre calendrier du Conseil d'arrondissement du 5~octobre\,». & \pg{01}{23} \\
11 sept.\ 2026 & Dépôt de la présente demande d'accès. & --- \\
11 et 14 sept.\ 2026 & Rencontre de travail n\textsuperscript{o}~2 (stationnement et hiver) et commentaires. & \pg{01}{167} \\
17 sept.\ 2026 & Livraison de la note préliminaire, filigrane «\,BROUILLON\,». & \pg{01}{167} \\
2 oct.\ 2026 & Note finale due. Aucun bordereau de transmission de la version finale ne figure dans les documents communiqués. & \pgl{01}{12}{12, 26} \\
6 oct.\ 2026 & Séance du conseil\,; annonce de la conversion hivernale. & --- \\
8 oct.\ 2026 & Dépôt d'une seconde demande d'accès (note finale et diapositives). & --- \\
9 oct.\ 2026 & Réponse à la présente demande. & --- \\
\bottomrule
\end{longtable}
\end{small}

\section{Constats détaillés}

\subsection{L'objectif était fixé avant toute analyse}

Le courriel du 10~juin, du chef de la Division des études techniques de
l'arrondissement à la Division de l'aménagement des rues de la Ville, est le
plus ancien document communiqué et le plus important \pgl{01}{7}{7--8}~:

\quot{«\,Les élus de l'arrondissement souhaitent revoir l'aménagement cyclable
de la rue Terrebonne réalisé en 2024, avec l'objectif de redonner la bande de
stationnement qui avait été retranchée. \ldots\ Les préoccupations
exprimées concernent principalement la largeur des aménagements cyclables,
jugée disproportionnée par rapport au volume de cyclistes observé. Selon
les élus, il serait possible de réduire l'espace consacré aux cyclistes
afin de rétablir la deuxième bande de stationnement retirée lors de la
reconfiguration de 2024.\,»}

\quot{«\,Vous comprendrez que nous avons déjà émis un avis technique quant
à la faisabilité de ces 2 concepts, toutefois, les élus souhaitent obtenir
un avis officiel, neutre et externe à l'arrondissement.\,»}

Trois choses en découlent. L'objectif, rétablir la bande de stationnement,
est présenté comme le but des élus, non comme une question. Les ingénieurs
de l'arrondissement avaient déjà donné un avis, et l'avis externe a été
demandé parce que les élus en voulaient un autre\,; la phrase ne dit pas
quel était l'avis interne, mais la description que fait le courriel des deux
concepts («\,une voie de circulation très étroite (moins de 4~mètres)
coincée entre deux bandes de stationnement\,»\,; «\,un aménagement
bidirectionnel très étroit\,» de 2,4~m avec une zone tampon de 0,3~m) et sa
prévision que le Service de sécurité incendie «\,pourrait s'y opposer\,»
laissent peu de doute. Enfin, la justification avancée pour l'exercice est
celle qui a servi en public depuis~: la piste est trop large pour le nombre
de cyclistes.

Le courriel se termine par un post-scriptum qui concerne les chiffres
utilisés par la suite~: «\,Nous disposons également d'un compteur permanent
Éco-Compteur sur le corridor. Au besoin, nous pourrons vous transmettre les
données disponibles\,» \pg{01}{8}. En juin 2026, les ingénieurs de
l'arrondissement savaient que le compteur permanent existait, le tenaient
pour la donnée pertinente et en offraient les données à la Ville.

L'«\,avis technique\,» interne n'était pas visé par ma demande et n'a pas
encore été demandé. Il devrait l'être.

Le document suivant dans la séquence est la directive du cabinet de la
mairesse du 13~août \pgl{01}{1}{1--2}, de la cheffe de cabinet au directeur
de l'aménagement urbain, avec copie au directeur d'arrondissement et aux
deux élus~:

\quot{«\,Après discussion, Mme la mairesse souhaite qu'une étude technique
formelle soit réalisée par une firme externe afin d'évaluer la
reconfiguration de la piste cyclable sur la rue Terrebonne. Le mandat serait
le suivant~: évaluer le scénario d'implantation d'une piste cyclable
bidirectionnelle d'un côté de la rue Terrebonne (entre le campus Loyola
et Côte-Saint-Luc) qui respecte les normes minimales prévues par la loi
et le retour du stationnement d'un côté de la rue.\,»}

Deux choses ressortent. Le mandat, tel que le cabinet l'a formulé, est
permanent et couvre toute la longueur~: une piste bidirectionnelle d'un
côté et le stationnement rétabli de l'autre, sans mention de l'hiver. Et
«\,qui respecte les normes minimales prévues par la loi\,» présuppose que les
concepts examinés en juin ne les respectaient pas~: les concepts décrits le
10~juin étaient une voie de circulation de moins de 4~m entre deux bandes de
stationnement et une piste bidirectionnelle de 2,4~m, les ingénieurs de
l'arrondissement avaient déjà donné leur avis à leur sujet, et le SUM y
avait vu «\,plusieurs compromis importants\,». La directive d'août demande
une version qui franchisse le minimum légal. Le retrait hivernal en est
absent\,; il apparaît une semaine plus tard dans la demande de services
rédigée par le personnel, à titre de mesure transitoire avant la mesure
permanente. Le 6~octobre, le conseiller Teodoresco a dit de cette
option~: «\,we try to keep bidirectional one side, putting it on one side
of the road, right? It was too complicated and no real upside for
anyone\,». En conversation avec l'auteur, il a dit que les ingénieurs de
l'arrondissement n'endosseraient pas une piste bidirectionnelle sur
Terrebonne, si légale soit-elle\,; le dossier est compatible avec les deux
affirmations, et la question qu'elles laissent ouverte est de savoir si le
Volet~2 demeure le plan.

\subsection{L'équipe vélo de la Ville et le Service de sécurité incendie}

L'analyse du SUM du 17~juin \pgl{01}{6}{6--7} est refusée en vertu des
articles~37 et 39. Sa conclusion survit dans les mots de l'arrondissement,
dans la demande du 22~juin au Service de sécurité incendie \pg{01}{6}~:

\quot{«\,À la demande de l'arrondissement, le Service de l'urbanisme et de la
mobilité a procédé à une analyse technique préliminaire des scénarios
envisagés. Cette analyse conclut que les options étudiées présentent
plusieurs compromis importants sur le plan de la sécurité, du confort des
usagers vulnérables et de la cohérence du réseau cyclable par rapport à
l'aménagement actuellement en place.\,»}

Le courriel d'accompagnement du SUM demande à être tenu informé «\,des
orientations qui seront retenues à la lumière de ces éléments\,»
\pg{01}{7}\,; les documents communiqués ne contiennent aucune réponse. La
réponse du Service de sécurité incendie du 13~juillet \pg{01}{5} se résume à
une phrase d'introduction suivie d'un caviardage qui couvre le reste de la
page. Quoi qu'elle dise, l'arrondissement a posé la question parce qu'il
s'attendait à une objection à une voie de moins de 4~mètres entre deux
rangées de voitures stationnées. Le 6~octobre, le conseiller a dit à la
séance que «\,we can keep the same capacity, the same level of safety for
our kids, and bring back \ldots\ more parking\,»\,; l'analyse du SUM est la
seule évaluation de sécurité du dossier dont la conclusion soit connue, et
c'est celle citée ci-dessus.

La séquence importe~: l'unité experte de la Ville a signalé des enjeux de
sécurité le 17~juin\,; puis l'arrondissement a retenu un consultant
externe le 20~août avec un mandat dont l'objet déclaré était d'augmenter le
stationnement. Une note rédigée selon ce mandat n'est pas une évaluation
indépendante de l'opportunité de retirer la piste\,; c'est un exercice de
conception portant sur la façon de le faire.

\subsection{Le mandat et ce qu'il demandait}

La demande de services du 20~août \pgl{01}{3}{3--4} est explicite quant à
l'objet et à la portée. Le titre~: une note technique pour évaluer une
piste bidirectionnelle d'un côté «\,et voir à évaluer la place du
stationnement sur rue, dans le but d'en augmenter le nombre\,». Parmi les
tâches~:

\quot{«\,Énumérer les mesures à mettre en oeuvre, en deux volets, pour le
retrait des aménagements cyclables durant la saison hivernale.
Volet~1 -- existant~: pistes cyclables unidirectionnelles bilatérales
(mesure transitoire pour l'hiver 2026--2027). Volet~1 alternatif~: maintien
et entretien hivernal d'une des deux bandes cyclables, et retrait de
l'autre. Volet~2 -- faisabilité d'une piste cyclable bidirectionnelle
unilatérale.\,»}

Le «\,Volet~1 alternatif\,» est, mot pour mot, la mesure annoncée le
6~octobre. Il figure dans les instructions de l'arrondissement au consultant
six semaines avant que celui-ci ne fasse rapport. Et il est qualifié de
transitoire~: le plan à plus long terme, dans le même document, est le
Volet~2, une seule piste bidirectionnelle avec une bande de stationnement
rétablie de l'autre côté à l'année. On n'a pas dit au conseil que la mesure
hivernale est la première étape d'une reconfiguration permanente, mais le
mandat le dit.

L'offre d'EXP \pgl{01}{9}{9--13} reprend la portée et ajoute les scénarios
hivernaux qu'elle comparerait~:

\quot{«\,\ldots\ la fermeture complète de la piste et son remplacement par du
stationnement en hiver; la fermeture en période de déneigement seulement;
ou, dans le cas des bandes unidirectionnelles actuelles, la fermeture
partielle avec une chaussée désignée et du stationnement des deux côtés
de façon permanente ou seulement lors des périodes de déneigement
durant la saison hivernale.\,»}

La formule limitée aux opérations de déneigement, que le conseiller a
écartée à la séance, était donc l'un des scénarios que le consultant était
payé pour évaluer. Si la note la déconseille, et pour quels motifs, c'est
précisément ce que le document refusé montrerait.

Prix~: 14\,755~\$ d'honoraires plus 2\,951~\$ de contingence, 17\,706~\$
avant taxes \pg{01}{16}, aux taux horaires de l'entente-cadre. Livrables~:
deux courriels intermédiaires, une note préliminaire, une note finale et
une réunion de clôture. Aucun travail de terrain, aucun comptage, aucun
relevé d'occupation\,; le consultant a reçu la présentation de 2025 de
l'arrondissement et le tableau du compteur le 28~août et a travaillé à
partir de ceux-ci et de l'imagerie aérienne de Google, Street View étant
périmé \pg{01}{19}.

\subsection{Un calendrier dicté par la date du conseil}

Le courriel du 18~août qui ouvre le dossier \pgl{01}{21}{21--22}~:

\quot{«\,\ldots\ la nouvelle mairesse souhaite que nous lui présentions une note
technique concernant la faisabilité de réaliser une piste cyclable
bidirectionnelle, unilatérale, sur la rue de Terrebonne et voir, en plus de
cette faisabilité, dans quelle mesure cet aménagement peut augmenter le
nombre de cases de stationnement sur rue. \ldots\ sachez que l'important
pour nos élus est que votre proposition puisse se faire très rapidement,
ainsi que la réalisation de la note en soi. L'objectif de nos élus est de
pouvoir agir sur vos conclusions pour un horizon début octobre, lors d'un
conseil d'arrondissement.\,»}

Et le 10~septembre \pg{01}{23}~: «\,Mon gestionnaire m'a rappelé hier qu'il
est très important que la note technique préliminaire nous soit acheminée
au plus tard le 18~septembre prochain pour respecter notre calendrier du
Conseil d'arrondissement du 5~octobre prochain.\,» La version préliminaire
du consultant, le 17~septembre, porte un filigrane «\,BROUILLON\,» parce que
«\,le délai serré ne nous a pas permis de faire faire une relecture
attentive\,» \pg{01}{167}. La note finale était due le 2~octobre, un
vendredi\,; l'annonce a été faite le mardi suivant.

\subsection{D'où viennent les chiffres}

Dans son allocution, le conseiller a dit (en anglais)~: «\,if we looked at
the official data provided by the borough and we have counters everywhere,
we have more counters than before. So we know exactly how many people use
it, when and how. So we're able to see that in winter, the volume of
traffic diminishes by 99\%. We go from 700 during the good months to 8.
That's 8 passage[s], right? So that's 4 cyclists, right? Including me\,»,
ajoutant que l'arrondissement disposait maintenant de «\,data for 3 years
in a row\,».\footnote{Transcrit de l'enregistrement de la séance du
6~octobre 2026 diffusé par l'arrondissement,
\url{https://www.youtube.com/live/0k_KZf9uxTk}.} En réponse à une question
du public, il a renvoyé aux diapositives et donné un chiffre de
12~déplacements à l'heure de pointe en janvier et un chiffre de 8. Les
diapositives sont dans les documents communiqués \pgl{01}{51}{51--73}~: la
présentation d'analyse des comptages de février 2025, mise à jour au cours
de 2025. Son tableau des vélos («\,MAJ 2025\,») donne, pour les deux périodes
de pointe de trois heures combinées, 12~vélos à Cavendish et 12 à Grand
toutes directions confondues en janvier 2025, et 10 et 8 sur l'axe
Terrebonne, contre 347 et 425 en juin 2025, étiquetés $-95$~\% et
$-96$~\% \pgl{01}{61}{61--62}. Les tableaux Miovision sous-jacents montrent
que les chiffres de janvier proviennent d'une seule journée, le 29~janvier
2025 \pgl{01}{85}{85, 155}\,; les totaux complets sur douze heures ce
jour-là étaient de 18 à Cavendish et de 12 à Grand sur l'axe Terrebonne.

Trois observations. Premièrement, il s'agit de comptages par caméra d'une
seule journée à deux intersections, lors du premier hiver de la piste,
effectués avant que le compteur permanent n'existe\,; le titre de la
diapositive, \emph{aux heures de pointes}, signifie que le huit est le
total des deux périodes de pointe de trois heures (6\,h\,30 à 9\,h\,30 et
15\,h\,30 à 18\,h\,30 \pg{01}{57})\,; ce n'est donc ni un chiffre
quotidien ni un chiffre horaire, et la diapositive elle-même ne donne
aucune indication qu'il s'agit du comptage d'une seule journée\,; la
lecture de 8 passages comme 4 cyclistes suppose que chacun est passé
devant la caméra aux deux périodes de pointe, ce que les tableaux
n'enregistrent pas. Deuxièmement, l'arrondissement disposait du relevé du
compteur permanent au moment de l'annonce~: la même liasse contient les totaux quotidiens du
26~juillet 2025 au 31~août 2026 \pgl{01}{74}{74--81}\,; les ingénieurs de
l'arrondissement avaient offert les données du compteur à la Ville en juin
\pg{01}{8} et envoyé le tableau au consultant le 28~août \pg{01}{19}. Il
montre 4\,891 passages entre le 1\textsuperscript{er}~décembre et le 31~mars
(41 par jour\,; 46 les jours de semaine\,; 63 les jours de semaine sur la
saison de déneigement du 16~novembre au 31~mars) et un deuxième été à
environ 700 par jour. Troisièmement, le «\,700\,» et le «\,8\,» que le
conseiller a opposés l'un à l'autre pour une baisse de «\,99~\%\,» ne sont pas
des chiffres comparables. Le 8, c'est Grand, axe Terrebonne seulement, six
heures de pointe, 29~janvier 2025 \pg{01}{62}. Le seul nombre voisin de 700
sur les diapositives est 707~: Grand, 19~septembre 2024, six heures de
pointe, mais sur \emph{les quatre approches}, y compris les vélos
circulant sur Grand elle-même, que le 8 exclut \pg{01}{61}\,; c'est aussi le
plus grand comptage pour une seule intersection de toute la présentation.
Diviser le plus petit nombre d'une diapositive par le plus grand de l'autre
donne 98,9~\%. Comparée de façon équivalente sur les diapositives
elles-mêmes, même intersection, même axe, mêmes six heures, la baisse de
septembre 2024 à janvier 2025 est de 98,2~\% (de 452 à 8) et, depuis les
comptages d'été de part et d'autre de cet hiver, d'environ 97~\% (de 283 et
281 à 8)\,; les diapositives elles-mêmes étiquettent janvier à $-95$~\%. Le
compteur permanent donne environ 93~\% entre la totalité de l'été et la
totalité de l'hiver. La baisse est forte selon toutes les mesures\,; le
«\,99~\%\,» est le seul appariement des nombres disponibles qui l'atteint.
Quant aux «\,data for 3 years in a row\,», les comptages de vélos
communiqués vont de septembre 2024 à août 2026 et couvrent deux hivers~: le
premier par la caméra, une seule journée\,; le second par le compteur
permanent.

La journée elle-même. Le 29~janvier 2025 était un mercredi pour lequel
Environnement et Changement climatique Canada avait émis un avis de temps
hivernal pour Montréal~: de cinq à dix centimètres de neige avec poudrerie,
les pires conditions devant durer jusqu'à midi, et la population invitée à
reporter les déplacements non essentiels.\footnote{Jessica Barile,
«\,Montrealers urged to delay non-essential travel as snowstorm hits, up to
10 cm possible\,», CityNews Montreal, 29~janvier 2025, 9~h~13,
\url{https://montreal.citynews.ca/2025/01/29/montrealers-urged-to-delay-non-essential-travel-as-snowstorm-hits-up-to-10-cm-possible/}.}
Les tableaux Miovision de Grand pour cette journée consignent tous les
modes, et non seulement les vélos \pgl{01}{82}{82--88}. Pendant les mêmes
six heures de pointe, la caméra a compté 2~298 automobiles sur les quatre
approches (617 sur les deux approches Terrebonne), 108 autobus, 20 camions
et 605 piétons, contre les 8 vélos sur l'axe Terrebonne, dont trois à la
pointe du matin et cinq à celle de l'après-midi. Par rapport au comptage
précédent à la même intersection, le 20~novembre 2024
\pgl{01}{131}{131--137}, les automobiles sur les approches Terrebonne
étaient en baisse de 23~\% (de 802 à 617), les piétons de 26~\% (de 818 à
605) et les vélos sur l'axe de 95~\% (de 154 à 8)\,; par rapport au
19~septembre 2024, les automobiles étaient en baisse de 11~\%, les piétons
de 31~\% et les vélos de 98~\%. Tous les modes ont reculé ce jour-là\,; le
vélo a reculé le plus, comme lors de toute journée de neige, et c'est
précisément pourquoi une seule journée ne mesure pas une saison. La
diapositive d'où le 8 a été lu n'indique ni la date ni la météo. Les
comptages par mode pour chaque journée Miovision du dossier sont tabulés
dans
\doc{2026-581/derived/miovision_counts_by_mode_2026-581.xlsx}{derived/miovision\_counts\_by\_mode\_2026-581.xlsx}.

Pour être complet~: la présentation ne contient aucune donnée de
Kensington, et ne le pouvait pas, la boucle ayant été mise en service le
26~juillet 2025, après la dernière mise à jour de la présentation (juin
2025). Sa diapositive de méthodologie mentionne bien l'Éco-Compteur,
«\,utilisé par le SUM depuis 2013\,», comme l'une des trois sources de
comptage avec Miovision et Telraam \pgl{01}{53}{29, 53}\,; les seules
données continues de la présentation sont celles des capteurs citoyens
Telraam de King-Edward et de Royal, qui donnent environ 242 vélos par jour
\pg{01}{65}. Le tableau du compteur est parvenu au consultant en pièce
jointe distincte du même courriel du 28~août que la présentation
\pg{01}{19}. Le consultant avait donc les deux documents dès la première
semaine du mandat, et l'arrondissement savait depuis juin au moins laquelle
des deux sources était la meilleure \pg{01}{8}\,; laquelle la note retient
ne pourra être connu qu'à sa communication.

\subsection{Ce que le dossier contient sur le stationnement}

L'objet déclaré de l'exercice est d'augmenter le stationnement\,; il vaut
donc la peine de noter ce que les documents communiqués contiennent de ce
côté. Le courriel du 10~juin énonce l'objectif et donne un seul chiffre,
géométrique~: la voie de circulation de moins de 4~mètres entre deux bandes
de stationnement \pg{01}{8}\,; rien sur la demande. La demande de services
charge EXP de «\,mesurer \ldots\ le nombre de cases de stationnement que cet
aménagement pourrait permettre de recréer\,» et de mesurer
«\,l'amélioration de la disponibilité du stationnement sur de
Terrebonne\,» \pgl{01}{3}{3--4}~: des questions posées au consultant, non
des réponses. L'offre d'EXP promet «\,une analyse de l'impact sur le nombre
et l'accessibilité des espaces de stationnement\,» et un livrable
«\,incluant l'impact sur le nombre de places de stationnement sur rue\,»\,;
sa liste de services ne comprend aucun relevé d'occupation ni aucun
comptage d'aucune sorte \pgl{01}{10}{10--11}. Le courriel du 28~août montre
que la première analyse intermédiaire portait sur le côté de la rue qui
devrait recevoir une piste bidirectionnelle, l'analyse du stationnement
devant suivre, et que le consultant travaillait à partir d'imagerie
aérienne \pg{01}{19}\,; le formulaire de démarrage prévoit l'«\,Analyse de
côté de rue du stationnement et aménagement hivernal\,» pour le
11~septembre \pg{01}{26}. La présentation des comptages ne contient aucune
diapositive sur le stationnement. La réponse au point~C2 de la demande, qui
visait la collecte de données derrière l'étude «\,including locations,
dates, seasons and hours of counts\,», est «\,the attached documents include
presentations\,», et les présentations sont des comptages de vélos et de
véhicules aux intersections.

Les documents communiqués ne contiennent donc aucun chiffre mesuré sur le
stationnement, de quelque nature que ce soit~: ni inventaire, ni
occupation, ni sondage, ni estimation des cases créées. Les «\,50 à 60
cases\,» que le conseiller a données à l'APC-NDG lors de la rencontre
qu'il a ensuite présentée au conseil comme sa consultation de l'Association
n'y figurent pas. Ce que la note refusée dit du
stationnement sera un décompte géométrique de cases qu'on pourrait marquer,
ce que le mandat demandait\,; rien dans le mandat ni dans les livrables ne
produirait une mesure du besoin, ni du nombre de nuits où ces cases
seraient nécessaires.

\subsection{Ce qui a été refusé, et pour quels motifs}

\begin{itemize}[nosep]
  \item \textbf{Article 39} (une analyse produite à l'occasion d'une
    recommandation, «\,jusqu'à ce que la recommandation ait fait l'objet
    d'une décision\,»)~: la note d'EXP dans toutes ses versions, l'analyse du
    SUM et les commentaires échangés sur les livrables.
  \item \textbf{Article 37} (avis ou recommandations faits depuis moins de
    dix ans, y compris par un consultant)~: les mêmes documents, plus l'avis
    du SIM.
  \item \textbf{Article 9, deuxième alinéa} (les ébauches, brouillons et
    notes préparatoires ne sont pas des documents de l'organisme public)~:
    les livrables intermédiaires et la note préliminaire.
  \item \textbf{Articles 23 et 24} (renseignements commerciaux de tiers)~:
    des parties des documents d'appel d'offres et de l'offre d'EXP\,; la
    ventilation des heures et des honoraires jointe à l'offre ne figure pas
    dans les documents communiqués.
  \item \textbf{Articles 53 à 59}~: renseignements personnels, appliqués aux
    noms et coordonnées dans les courriels\,; en pratique, les noms du
    personnel restent visibles.
\end{itemize}

L'article~39 est le motif qui porte tout le poids, et il est limité dans le
temps par ses propres termes. Le conseiller a annoncé la mesure le
6~octobre, en disant que les élus étaient «\,now ready to give our
orientation to the services, and the services will tell you how the
mechanics of what the adjustment will happen during this winter\,» et que
le retrait saisonnier «\,is what we're going to propose. It's going to be a
test\,». Sa déclaration à la \emph{Gazette} du 7~octobre dit que «\,one
thing is certain: residents will see changes by Christmas 2026\,», tout en
disant que le changement «\,still need[s] to go through several steps\,». Une
déclaration signée, publiée sur sa page de médias sociaux dans les jours
suivant la séance, réunit les deux moitiés dans ses propres mots~:
«\,The solution we are currently developing with experts is to maintain the
bike lane as is during the summer and remove it on one side only during the
winter season. Several steps remain before this proposal can be
implemented, including consideration by the city's decision-making bodies,
namely the borough council and City Council. One thing is certain:
residents will see changes before Christmas 2026.\,» (La solution que nous
élaborons actuellement avec des experts consiste à maintenir la piste telle
quelle l'été et à la retirer d'un seul côté l'hiver. Plusieurs étapes
restent à franchir, dont l'examen par le conseil d'arrondissement et le
conseil municipal. Une chose est certaine~: les résidents verront des
changements avant Noël 2026.)
L'arrondissement ne peut pas avoir les deux à la fois~: si la recommandation
a fait l'objet d'une décision, l'article~39 ne s'applique plus et la note
doit être communiquée sous la seule réserve de l'article~37, qui est
discrétionnaire et ne couvre pas le contenu factuel\,; si elle n'en a pas
fait l'objet, l'annonce publique d'un résultat décidé était prématurée.
L'article~14 de la Loi exige que le contenu factuel soit extrait et
communiqué lorsqu'une restriction ne vise qu'une partie d'un document\,; la
réponse ne le fait pas pour la note.

Le motif de l'article~9 est faible pour la version préliminaire du
17~septembre. C'était un livrable contractuel, transmis au client pour
révision, et les commentaires de l'arrondissement ont façonné la version
finale. Le mot «\,brouillon\,» en filigrane n'en fait pas une ébauche
personnelle au sens de l'article~9, et la Commission l'a dit dans des cas
comparables.

\section{Ce que le dossier établit et n'établit pas}

\subsection{Sur le processus décisionnel}

Le dossier montre un objectif fixé par les élus au printemps\,; un avis
technique interne déjà donné par les ingénieurs de l'arrondissement, après
quoi un avis externe a été sollicité\,; une analyse de sécurité de l'équipe
vélo de la Ville reçue en juin et jamais mentionnée en public\,; un mandat
rédigé en août dont l'objet déclaré était d'augmenter le stationnement,
avec une date de conseil déjà choisie\,; et une présentation de comptages
qui a écarté les données du compteur permanent de l'arrondissement au
profit d'un comptage d'une seule journée de l'hiver précédent, et aucune
mesure de la demande de stationnement. Que les élus
aient raison ou non de vouloir plus de stationnement est une autre
question\,; ce que les documents établissent, c'est l'ordre dans lequel la
décision et sa justification ont été produites, et cet ordre diffère du
récit donné à la séance du 6~octobre.

\subsection{Sur l'option limitée aux opérations de déneigement}

L'arrondissement a payé pour une évaluation qui comprenait «\,la fermeture
en période de déneigement seulement\,» et a dit à l'APC-NDG, lors de la
rencontre qu'il a ensuite citée au conseil comme sa consultation de
l'Association, que l'option avait été écartée parce qu'on ne pouvait en
informer les résidents\,; l'option n'a pas été mentionnée dans l'allocution
du 6~octobre. Le raisonnement, s'il va plus loin, se trouve dans la note
refusée.

\subsection{Sur ce qui est prévu au-delà de cet hiver}

Le cabinet de la mairesse a demandé le 13~août une piste bidirectionnelle
permanente avec retour du stationnement\,; «\,mesure transitoire pour l'hiver
2026--2027\,» et Volet~2, ensemble, décrivent la conversion hivernale comme
une première étape vers celle-ci. C'est une mesure différente de celle
qui a été annoncée, avec un profil de sécurité différent~: les «\,compromis
importants\,» du SUM concernaient précisément ces coupes de rue. Elle n'a
pas été divulguée au public. La déclaration du conseiller selon laquelle la
piste sera maintenue «\,as is during the summer\,» est le premier propos
public qui s'y rapporte\,; elle ne dit pas si l'option bidirectionnelle a
été abandonnée, et la question devrait être posée directement. Dans son
allocution, le conseiller a décrit l'option bidirectionnelle comme «\,too
complicated and no real upside for anyone\,», le retrait hivernal comme
«\,a test\,» («\,We're going to see the data this year, and hopefully try to
make it even better after\,») et l'ensemble comme «\,part of a longer term,
more sophisticated bike plan \ldots\ for next year\,». La mesure porte donc
trois descriptions au dossier~: «\,mesure transitoire pour l'hiver
2026--2027\,» dans la demande de services, «\,seasonal\,» et «\,a test\,»
dans l'allocution. La même déclaration nomme le conseil municipal parmi les instances qui doivent
examiner la proposition, ce qui, si c'est le cas, place le service
d'aménagement des rues de la Ville elle-même, auteur de l'analyse du
17~juin, sur le chemin de l'approbation.

\subsection{Au regard de la méthode de révision de la Ville}

\textbf{Deux mandats.} La piste elle-même a été construite à la suite
d'une étude d'une autre nature. L'administration élue en 2021 avait été
transparente, dans sa campagne, sur son engagement à étendre le réseau
cyclable, et l'étude EXP de 2023 qu'elle a commandée sur Terrebonne n'était
que la diligence raisonnable en cours de route vers la réalisation de cet
engagement~: une étude de faisabilité et de conception pour un projet déjà
décidé, dont la tâche était de confirmer l'absence d'obstacle et d'arrêter
l'aménagement, et dont les limites, que les opposants à la piste ont
relevées à l'époque, étaient celles d'une étude de conception. En matière
d'infrastructures cyclables, la campagne de l'administration élue en 2025
reposait sur un engagement à fonder les décisions sur des études~: non pas
sur le retrait de pistes, mais sur «\,un audit complet de
toutes les infrastructures cyclables de Montréal dans les 100 premiers
jours\,», afin de «\,revoir, sécuriser ou retirer les infrastructures
cyclables là où c'est nécessaire\,», la cheffe se décrivant comme
«\,pro-vélo\,». Le conseiller a lui-même décrit cet engagement, le
6~octobre, comme pris au porte-à-porte~: «\,when you do door knocking on
1000s of doors, and you see the most important issue for people there
\ldots\ the bike path that has become such a synonym of not being
consulted \ldots\ we have made a commitment and we have to act on this
promise\,».\footnote{CBC News, «\,How do Montreal's top mayoral
candidates compare on election issues?\,» (comparaison des plateformes,
octobre 2025),
\url{https://newsinteractives.cbc.ca/features/2025/montreal-municipal-election/}\,;
CBC News, «\,Are Montreal bike paths a boon or burden? See what mayoral
hopefuls are promising\,», \url{https://www.cbc.ca/1.7648562}\,; La
Presse, «\,Martinez Ferrada veut revoir l'ensemble du réseau
cyclable\,», 11~août 2025,
\url{https://www.lapresse.ca/actualites/grand-montreal/2025-08-11/elections-municipales/martinez-ferrada-veut-revoir-l-ensemble-du-reseau-cyclable.php}.
Les citations de la plateforme sont traduites de l'anglais de CBC.} Une
étude appelée à décider si une piste reste porte un poids qu'une étude de
conception ne porte pas, et l'administration a énoncé le 10~février 2026
ce qu'une telle étude comporte.

La norme que l'administration s'est donnée pour réviser une piste cyclable
est publique. Le communiqué de la Ville du 10~février 2026 lançant son audit
du réseau cyclable décrit la méthode comme «\,recueillir les données,
analyser et corriger au besoin\,», avec «\,analyse technique, collecte de
données et observation terrain\,» sur quatre pistes existantes (Rachel,
Côte-Sainte-Catherine, de la Commune, Meilleur) pendant la haute saison,
sous la direction du Service de l'urbanisme et de la mobilité, des résultats
attendus vers la fin de 2026, et une révision du Plan vélo pour établir
«\,des critères clairs, cohérents et harmonisés\,»\,; la mairesse a expliqué
que «\,les documents qu'on nous a remis ne répondaient pas aux préoccupations
exprimées sur le terrain\,», d'où le choix de recueillir d'abord ses propres
données. Le même communiqué nomme Terrebonne, avec Hochelaga, Gouin et
Henri-Bourassa, parmi les rues où la Ville travaille avec l'arrondissement
sur la sécurité.\footnote{Ville de Montréal, «\,10 actions en 100 jours~: la
Ville de Montréal annonce un état des lieux pour une mobilité
planifiée\,», 10~février 2026,
\url{https://montreal.ca/communiques-de-presse/10-actions-en-100-jours-la-ville-de-montreal-annonce-un-etat-des-lieux-pour-une-mobilite-planifiee}\,;
Journal Métro, «\,Montréal lance son audit des pistes cyclables\,»,
10~février 2026,
\url{https://journalmetro.com/actualites/politique/3208452/montreal-lance-son-audit-des-pistes-cyclables/}.}
Ce processus est toujours en cours. Outre l'audit, la Ville a tenu une
consultation en ligne sur le réseau cyclable en juin et juillet 2026\,;
son propre échéancier situe l'analyse des résultats et la poursuite de la
consultation publique («\,détails à venir\,») à l'automne 2026, la
rétroaction auprès de la population au printemps-été 2027 et le nouveau
plan vélo à l'automne 2027.\footnote{Ville de Montréal, Réalisons
Montréal, «\,La place du vélo à Montréal\,», échéancier du projet,
\url{https://realisons.montreal.ca/projects/place-du-velo} (consulté le 13~octobre 2026)\,; CBC News, «\,Montreal cyclists
question objectivity, privacy in city's survey on bike lanes\,»,
1\textsuperscript{er}~juillet 2026.} Mesuré à cette méthode, le dossier
Terrebonne montre un objectif fixé avant l'analyse, une analyse du SUM
défavorable et refusée, un mandat dont l'objet était le stationnement,
aucune collecte de données d'aucune sorte, et une annonce faite alors que
l'audit, la consultation et le plan de la Ville étaient toujours en cours.
Que les deux processus soient formellement liés est une question pour
l'administration\,; le contraste de méthode et l'ordre des dates, eux,
sont des affaires de dossier.

\subsection{Ce que le dossier n'établit pas}

La conclusion du SUM n'est connue que telle que l'arrondissement l'a
résumée\,; l'avis du SIM est inconnu\,; le contenu de la note d'EXP est
inconnu, et il se peut fort bien qu'elle déconseille les options plus
étroites et recommande la mesure hivernale comme la moins mauvaise.
«\,Disproportionnée\,» est le compte rendu, par le personnel, de l'opinion des
élus, non un constat. Aucune des citations ci-dessus ne devrait être
poussée au-delà de ce qu'elle dit. Elles établissent la séquence et
l'objet\,; elles n'établissent pas ce que le consultant a conclu, ce qui est
la raison d'obtenir la note.

\section{Index des documents}

Les numéros de page renvoient à la page du PDF de chaque liasse
communiquée.

\begin{small}
\begin{longtable}{>{\raggedright\arraybackslash}p{3.0cm}>{\raggedright\arraybackslash}p{9.2cm}>{\raggedright\arraybackslash}p{2.6cm}}
\toprule
Document & Contenu & Emplacement \\
\midrule
\endhead
Lettre de réponse & Réponses point par point\,; motifs de refus\,; avis de recours. & \doc{2026-581/raw/response_2026-10-09.pdf}{Request 2026-581}, 8~p. \\
Liasse 01 & & \doc{2026-581/raw/01_emails_mandate_counts.pdf}{01 Doc\_req}, 167~p. \\
\quad Directive du cabinet & Courriel du 13~août de la cheffe de cabinet fixant le mandat (piste bidirectionnelle, retour du stationnement), transmis les 13 et 14~août. & p.~1--2 \\
\quad Chaîne SIM / SUM & Demande au SUM du 10~juin\,; réponses des 12 et 17~juin\,; demande au SIM du 22~juin\,; réponse du SIM du 13~juillet (caviardée). & p.~5--8 \\
\quad Demande de services & 20~août 2026, mandat 59-26-ET-01\_002. & p.~3--4 \\
\quad Offre d'EXP & 25~août 2026\,; portée, scénarios, échéancier, 14\,755~\$. & p.~9--13 \\
\quad Bon de commande & 26~août 2026, n\textsuperscript{o}~1796945, 17\,706~\$. & p.~15--17 \\
\quad Courriels avec EXP & 18 au 28~août~: mandat, plan de travail, intrants, première analyse intermédiaire. & p.~19--22 \\
\quad Démarrage & Formulaire de la réunion du 3~septembre, livrables et dates\,; courriels du 10~septembre sur l'échéance du conseil du 5~octobre. & p.~23--26 \\
\quad Présentation des comptages & «\,Suivi de projet -- Analyse des comptages\,», février 2025, et mise à jour «\,MAJ 2025\,». & p.~27--49\,; 51--73 \\
\quad Boucle de comptage & Totaux quotidiens Kensington, vélos et véhicules motorisés, 26~juillet 2025 -- 31~août 2026. & p.~74--81 \\
\quad Tableaux Miovision & Comptages aux intersections par période de 15~minutes et par classe~: Grand, Cavendish, Walkley, Royal\,; juillet 2024 -- septembre 2025\,; le 29~janvier 2025 aux p.~82--88 et 152--158. & p.~82--164 \\
\quad Livraison de la version préliminaire & Transmission de la note préliminaire, 17~septembre 2026. & p.~167 \\
Liasse 02 & Appel d'offres pour l'entente-cadre 59-25-ET-01, règlement, dossier décisionnel (p.~169--177), grille d'évaluation, soumission de WSP. & \doc{2026-581/raw/02_framework_agreement_tender.pdf}{02 DOC requis}, 341~p. \\
Résolution & CA25~170126, 9~juin 2025~: ententes-cadres à WSP (1,60~M\$) et à EXP (1,70~M\$). & \doc{2026-581/raw/03_resolution_CA25-170126.pdf}{03 Résolution}, 2~p. \\
Extraits de la Loi & art.~9, 23--24, 37, 39, 53--59\,; avis de recours. & \href{\repourl/tree/main/2026-581/raw/act_extracts}{act\_extracts/}, 5~fichiers \\
Dérivé & Totaux quotidiens Kensington extraits en CSV à partir des p.~74--81. & \doc{2026-581/derived/kensington_daily_borough_release_2025-07_2026-08.csv}{CSV} \\
Dérivé & Comptages Miovision par mode (automobiles, autobus, camions, vélos, piétons), six heures de pointe et 12~h, toutes approches et axe Terrebonne, pour chaque journée de comptage, d'après les p.~82--164\,; le 29~janvier 2025 à Grand détaillé. & \doc{2026-581/derived/miovision_counts_by_mode_2026-581.xlsx}{XLSX}, \doc{2026-581/derived/miovision_counts_by_mode_2026-581.csv}{CSV} \\
\bottomrule
\end{longtable}
\end{small}

\end{document}
